\section{Evaluation}
\subsection{Questions and experimental contract}
We ask whether the separated boundaries help a complete scientific analysis,
whether fixed concurrency or dependency initialization explains the gain,
and how the result changes across physical hosts. The primary evidence comprises
\UpgradeTrials{} accepted trials in five randomized blocks per configuration.
Figures show medians and observed ranges. With five samples we report paired
ratios and wins, not precise population confidence intervals.

Colocated controls use two four-core workers with disjoint affinity within an
eight-CPU allocation. Application studies use their own
allocation. The control continuation retains complete comparison blocks and
reruns an interrupted block in full; it never pairs runtimes from different
allocations. Physical-host scaling holds eight persistent batch allocations
on eight distinct hosts of the same machine class and selects nested subsets
of one, two, four, and eight hosts, with eight CPUs per worker. All task records
must match an admitted worker's declared CPU set. These are CPU shares with
explicit affinity, not exclusive whole-node reservations. Other users and
submission-host load can still introduce contention.

Each trial starts a fresh service or manager, admits its complete pool, then
times graph construction, submission, execution, every sink fetch, and local
sink fsync. DataVine also admits requested outputs durably at its controller;
the persistence operation count is consequently not identical. Worker-local
intermediates are used in both backends. Pool admission is excluded and
recorded separately. Libraries have no additional readiness barrier, so
residual library startup remains inside the interval.
The TaskVine baseline uses the ordinary FuturesExecutor and temporary-output
path in the same checkout, including its shared function-credit implementation;
it is not a separately built, unmodified upstream release.

Python is 3.10 with NumPy 1.26.4; the analysis additionally uses Uproot 5.6.8,
Awkward 2.8.7, and Vector 1.6.3. Dependencies are copied to node-local storage
before colocated campaigns. Both libraries explicitly preload the same
modules before forking, except in the labeled cold controls. DataVine's
generated executor exists before service startup, and its hash must match
the executable actually received by every worker. This check caught and
excluded an earlier harness version that silently deployed the default
executor. Failed, interrupted, and incompatible trials remain diagnostic
artifacts; none contributes a zero bar or a replacement success.

\subsection{Serialized scheduling and concurrent data-plane work}
The first architectural claim is that data-lifecycle service demand need not
consume the manager's scheduling reactor. In the conventional single-reactor
path, dependency scheduling and dispatch share that reactor with replica or
checkpoint coordination, stage-in/out acknowledgements, and retention or
garbage-collection updates. The workflow process retains one owner for
logical scheduling and TaskVine manager state, but the service endpoint handles
data-plane RPC connections with an independent event-thread pool, while
requested-output persistence runs in a separate Controller data-worker pool.
The current-build validation exposes both counts as
\texttt{DATAVINE\_RPC\_THREADS} and \texttt{DATAVINE\_DATA\_THREADS}; defaults
remain one RPC thread and 16 data workers. The manager never becomes the path
for the durable output bytes, although it still receives the corresponding
completion and lifecycle notifications.

Table~\ref{tab:controller-decoupling} holds the single-thread scheduling
service and workload fixed while changing only the Controller data-worker
count. Two repetitions at each point produced 26.9, 43.9, and 63.7 tasks/s
medians for one, four, and 16 data workers, respectively, a 2.37$\times$
median improvement from one to 16. Every run completed 256 one-MiB outputs
and verified all 268 MiB of durable bytes. The manager's owner-execute time
was 6.9--27.6 ms across these runs, while the end-to-end service interval was
3.1--14.6 s. This separates the manager's scheduling control path from the
measured durability work, while exposing local-filesystem variability. The
experiment does not claim that all data operations scale linearly; it tests
whether data-plane work that would otherwise share the manager's serialized
service path can be served concurrently.

\input{results/controller-decoupling-table.tex}

The worker-local control completes the same 256 tasks at 461.7 and 453.1
tasks/s with one and 16 data workers, respectively, and creates no durable
Controller files. Manager owner-execute time remains 1.59 and 1.58 ms. Thus
the thread-sensitive cost is in requested-output data management, while the
logical scheduling path remains stable.

An unbatched metadata-RPC control processed 39.5k publish and 37.8k resolve
records/s with one client connection. Sixteen concurrent connections reached
91.5k and 84.6k records/s even with one RPC event thread; 16 event threads
reached 98.9k publish records/s. The small incremental gain and the resolve
plateau are expected: metadata still uses shared controller state. The result
supports independent data-plane concurrency and identifies its remaining
serialization boundary; it does not claim unlimited Controller scaling.
The complete protocol regression with 16 RPC and 16 data threads passed
submission, restart, compare-and-swap append, journal recovery, corruption
rejection, the 1024-connection cap, and object-store deduplication.
The raw summary and source fingerprints are retained with the artifact.

\subsection{Complete ATLAS Open Data analysis}
We port the existing diphoton analysis accompanying ATLAS Open Data educational
materials~\cite{atlas-open-data}. The complete available sample contains 16
ROOT files, 9,861,498,743 bytes, and 36,564,144 collision events. The DAG has
374 read/decompression tasks, 374 photon-selection tasks, and 54 reduction
tasks: 802 individual attempts. Reads cover at most 100,000 events each;
selection consumes the actual jagged arrays, and an eight-way reduction
combines selection counts and the diphoton-mass histogram. There are no
synthetic sleeps or substituted numerical kernels in this application.

An independent run of the existing full-file analysis pins every input hash,
all seven selection counts, and all 60 histogram bins. Every accepted runtime
run agrees exactly, including 553,456 selected events. We preserve the reference
selection, including its permissive eta OR predicate; the purpose is runtime
equivalence, not a revised physics selection. A diagnostic implementation that
rearranged float32 mass arithmetic changed five selected events and was rejected.
The accepted version preserves the reference Vector operation order.
Exact numerical acceptance is platform-specific: a TaskVine replay on a second
platform selected 553,458 events against the primary reference's 553,456.
The unchanged independent reference on that platform reproduced every observed cut
and bin difference. We retain the rejected run, provide independently generated
platform-specific manifests, and do not assume cross-platform bitwise identity.
The primary performance figures continue to use the original allocation.

The complete ROOT dataset is copied and hashed once before the application
campaign. That setup cost is retained but excluded from workflow time. Reading,
decompression, selection, array serialization, intermediate delivery, reduction,
and durable sink admission remain inside the timer. OS caches are not flushed.
This is a complete educational analysis over real data on one compute host;
it is not a production collaboration analysis or distributed input-store study.

\begin{figure}[t]
\centering\includegraphics[width=\columnwidth]{upgrade-application.pdf}
\caption{The full 36.56-million-event ATLAS analysis. Five paired blocks,
802 individual tasks per run, and exact agreement with the independent oracle.
Bars are medians; whiskers are observed ranges.}
\label{fig:application}
\end{figure}
DataVine fixed and elastic complete in median \UpgradeAtlasFixedMedian{} and
\UpgradeAtlasElasticMedian{} seconds, versus \UpgradeAtlasTVMedian{} seconds
for TaskVine. Median paired completion-time ratios are
\UpgradeAtlasFixedSpeedup{} and \UpgradeAtlasElasticSpeedup{}, respectively.
The close fixed/elastic result is evidence against attributing this application
gain primarily to feedback. The comparison supports the implemented execution
path under the stated initialization and storage contract.

\subsection{Matched concurrency and preparation controls}
CPU chains contain 64 independent lanes of eight 50-ms process-CPU tasks,
with 1-KiB outputs. Hash chains use the same 512-task graph with 1-MiB outputs;
each consumer checks its parents' payload hashes. The phase DAG contains 256
tasks alternating 100-ms CPU work and actual node-local writes, fsync, and
reads with 256-KiB payloads and fan-in. Task bodies, graph structure, logical
result visibility, and sink acceptance are identical across backends.

\begin{figure*}[t]
\centering\includegraphics[width=.96\textwidth]{upgrade-controls.pdf}
\caption{Concurrency and preparation controls, five paired blocks each.
DV fixed uses one execution slot per core; DV elastic adapts its window;
DV eager also advances preparation. TV 4C is the explicitly extended fork
library with four slots per allocated core, not four times the CPU allocation.}
\label{fig:controls}
\end{figure*}
\begin{table}[t]
\centering\small
\caption{Median completion time in seconds for the matched controls.}
\input{results/upgrade-controls-table.tex}
\end{table}
The extended TaskVine control removes only the generated fork-library startup
guard against slots exceeding cores. Its physical allocation remains fixed.
It challenges the explanation that extra concurrency alone is sufficient.
Within DataVine, fixed, elastic, and eager controls are close on these graphs;
there is no consistent case here for making the adaptive window the principal
research contribution. End-to-end differences also include serialization,
graph submission, and result processing. They are not a factorial isolation
of scheduler and data-controller costs.

\subsection{Scaling across distinct physical hosts}
\begin{figure*}[t]
\centering\includegraphics[width=.96\textwidth]{upgrade-scaling.pdf}
\caption{Fixed-size 2,048-task graphs on nested subsets of eight persistent,
distinct physical hosts, eight CPU shares per host. Five paired blocks;
medians and ranges. Nodes are not exclusively reserved.}
\label{fig:scaling}
\end{figure*}
The scaling graph has 256 lanes and eight stages. CPU tasks consume 100 ms
each and emit 1 KiB; the data case generates and hashes 1 MiB per task with
minimal additional CPU work. Separating these cases avoids interpreting a
CPU scaling result as a data-path scaling result. The same persistent pool
is used throughout; a continuation retains all completed non-grouping trials
and resumes the original randomized order after incompatible grouping trials
are removed from the comparison.

\input{results/upgrade-scaling-summary.tex}
The data-path result is particularly consequential: more remote endpoints can
add scheduling and delivery work without shortening the critical path. These
measurements bound the useful envelope of this implementation. They do not
establish controller capacity at 16 or more nodes, nor isolate whether manager
serialization, placement, or data-controller service dominates every plateau.

\subsection{Initialization, instrumentation, and CPU accounting}
The attribution campaign holds a single allocation and executes five paired
blocks at 1-KiB and 1-MiB payload sizes. It compares preloaded and cold libraries,
and instrumented and uninstrumented execution, independently for each backend.
This avoids inferring a cold-start effect from two different hardware campaigns.

\input{results/upgrade-attribution-summary.tex}
The cold controls retain DataVine's startup disadvantage rather than treating
preloading as a free default. Users must provision an explicit, fork-compatible
dependency environment. The prototype's executable override is now fail-closed:
an explicitly missing or non-executable path cannot silently select a different
executor. Performance evidence remains attached to its exact measured binaries;
the subsequent compatibility repairs have separate regression evidence.

Worker-lifetime CPU accounting includes worker and library startup and shutdown.
Task-body CPU excludes deserialization and some import costs. Optional process
sampling is non-atomic, and summed RSS can count shared pages repeatedly.
\begin{table}[t]
\centering\small
\caption{Median CPU seconds. DV uses fixed admission. Worker lifetime aggregates
both worker process trees; it excludes manager, controller, and client CPU.}
\input{results/upgrade-cpu-table.tex}
\end{table}
Similar task-body CPU and different worker-lifetime CPU support an explanation
involving runtime and library overhead, rather than cheaper numerical work.
They do not establish total cluster CPU savings: some work can move to the
controller, and the two timer scopes must not be subtracted as a full accounting.
The artifact separately reports these scopes and manager send/status timers.
Manager file-byte counters exclude peer traffic; we do not label them total
network bytes. Validation of 5,120 local-clock traces confirms one ordered
trace per instrumented physical attempt. Reconstructed stage occupancy counts
describe tasks in each stage, not physical prepared-data bytes.

\subsection{Correctness, failure injection, and remaining scope}
The retained measured build passes \UpgradeRegressionPassed{} DataVine regression
gates. The current configurable service build additionally passes the native
workflow protocol regression with both one and 16 RPC event threads and one and
16 Controller data workers.
Explicit invalid-executor tests require failure before service contact.
The runner verifies the loaded Python extension and records binary hashes,
preventing an installed older extension from masquerading as this build.

The retained recovery experiment removes three owned worker jobs from a
1,024-task workflow and validates all identities and 16 sinks after 18
resubmissions. Focused eager/deferred runs validate 32 ordered attempts per
mode, and generic serverless execution remains covered. These are exercised
semantics, not exhaustive failure coverage. Whole-node 16+ scaling,
prepared-node assignment, and exact sender-level traffic and prepared-byte
accounting remain open. Requested-output fsync is not controller-host-loss
durability, and no arbitrary external-side-effect exactly-once claim is made.
